\documentclass[9pt,twoside]{pnas-new}

% Use the lineno option to display guide line numbers if required.
\usepackage{adjustbox}  % in preamble
\usepackage{booktabs}
\usepackage{multirow}
\usepackage{graphicx}
    
\usepackage{pdfpages}
\usepackage{amsmath}
\usepackage{amsthm}
\usepackage{amssymb}
\usepackage{setspace}
\usepackage{graphicx}
\usepackage{booktabs}
\usepackage{pdflscape}
\usepackage{url}
\usepackage{caption}
\usepackage{multirow,rotating,multicol}
\usepackage{listings}
\usepackage{longtable}
\usepackage{threeparttable}
%\usepackage[position=top, font=normalsize]{subfig}
\usepackage{subcaption}
\usepackage{float}
\usepackage{pgfpages,tikz}
\usepackage{adjustbox}
\usepackage{enumitem}
\usepackage{appendix}
\usepackage{dsfont}
\usepackage{natbib}
\usepackage{bookmark}

\templatetype{pnassupportinginfo}


\title{Beyond Stoves and Outdoor Monitors: Indoor PM2.5 Exposure and Inequality in a Fast Growing Megacity}

% Use letters for affiliations, numbers to show equal authorship (if applicable) and to indicate the corresponding author

% \title{Your main manuscript title}
\author{Sun, Sorin, Suryadarma, Rahut and Resosudarmo}
\correspondingauthor{Sorin \\E-mail: jsorin@uchicago.edu}

\begin{document}

\maketitle

\renewcommand{\thefigure}{S\arabic{figure}}
\renewcommand{\thetable}{S\arabic{table}}


%% Adds the main heading for the SI text. Comment out this line if you do not have any supporting information text.
% \SItext


% \subsection*{Subhead}
% Type or paste text here. This should be additional explanatory text such as an extended technical description of results, full details of mathematical models, etc.   

% \section*{Heading}
% \subsection*{Subhead}
% Type or paste text here. You may break this section up into subheads as needed (e.g., one section on ``Materials'' and one on ``Methods'').

% \subsection*{Materials}
% Add a materials subsection if you need to.

% \subsection*{Methods}
% Add a methods subsection if you need to.
\newpage


\SItext
\begin{adjustwidth}{0cm}{1cm} 
\section*{Survey Instruments}

\subsection*{Baseline Survey Instruments}

Study participants were asked the following survey questions during a baseline survey, before the installation of the air quality monitoring hardware.\footnote{These questions were part of a longer survey.} Bahasa Indonesian translations are in parentheses.

\begin{enumerate}
\item What was your household's total income in the past month (in IDR)? This includes cash transfers from the government (\textit{Berapa penghasilan rumah tangga dalam sebulan terakhir? (dalam Rupiah)})
\begin{enumerate}
    \item $<$Rp 1.000.000
    \item Rp 1.000.000 - Rp 2.999.999
    \item Rp 3.000.000 - Rp 4.999.999
    \item Rp 5.000.000-Rp 7.999.999
    \item Rp 8.000.000 - Rp 9.999.999
    \item Rp 10.000.000 - Rp12.999.999
    \item $>$13.000.000
        \item Does not know (\textit{Tidak tahu})
        \item Refuse to answer (\textit{Menolak menjawab})
\end{enumerate}
\item How many rooms does this residence have? (\textit{Berapa banyak kamar yang dimiliki hunian ini?})

[Enumerator: Count spaces as separate rooms only if they have a formal entrance, even if the door is a curtain. Do not count one room divided by curtains without a formal entrance as separate rooms. Family-owned shops attached to the residence and bathrooms should count as rooms.] \\ 
(\textit{[Enumerator: ruangan dianggap terpisah jika ada tempat masuk, meskipun pintunya adalah kain. Jika tidak ada tempat masuk, meski terpisah oleh kain/gorden, jangan dihitung sebagai ruangan terpisah. Toko yang bersatu dengan rumah, kamar mandi, dihitung sebagai satu ruangan terpisah.]})

    \item In the past 7 days, how do you think air quality was in your home, on the scale we mentioned earlier? (\textit{Dalam seminggu terakhir, berdasarkan skala yang sudah saya sebutkan tadi, menurut Ibu/Bapak seberapa baik kualitas udara di Jakarta?})

    \begin{enumerate}
        \item Good (green) (\textit{Baik (hijau)})
        \item Moderate (yellow) (\textit{Sedang (kuning)})
        \item Unhealthy for sensitive groups (orange) (\textit{Tidak sehat untuk kelompok rentan (oranye)})
        \item Unhealthy (red) (\textit{Tidak sehat (merah)})
        \item Very unhealthy (purple) (\textit{Sangat tidak sehat (ungu)})
        \item Hazardous (Brown) (\textit{Berbahaya (coklat)})
        \item Doesn't know (\textit{Tidak tahu})
    \end{enumerate}

    
    \item Is there an air conditioner in the house? (\textit{Apakah terdapat AC/pendingin udara di dalam rumah?})
    \begin{enumerate}
        \item No (\textit{Tidak})
        \item Yes, in room the hardware is in (\textit{Ya, di ruangan yang terdapat air purifier dan air monitor})
        \item Yes, in another room in the house (\textit{Ya, di ruangan lain})
    \end{enumerate}
     
    \item In the past 7 days, how many times did you see trash being burned in your neighborhood? (\textit{Dalam seminggu terakhir, berapa kali Ibu/Bapak melihat sampah dibakar di lingkungan Ibu/Bapak?})
    \begin{enumerate}
        \item Never (\textit{Tidak pernah})
        \item 1 or 2 times (\textit{1-2 kali})
        \item 3 or more times (\textit{3 kali atau lebih})
        \item Does not know (\textit{Tidak tahu})
        \item Refuse to answer (\textit{Menolak menjawab})
    \end{enumerate}
    \item What is the main cooking fuel you use in your home? (\textit{Apa bahan bakar utama untuk memasak yang Ibu/Bapak gunakan di rumah?})
    \begin{enumerate}
        \item Electric (\textit{Kompor listrik})
        \item Gas/LPG (\textit{LPG})
        \item Kerosene (\textit{Minyak tanah})
        \item Charcoal/briquettes/coconut shells (\textit{Arang/briket/batok kelapa})
        \item Firewood (\textit{Kayu bakar})
        \item Other (\textit{Lainnya})
    \end{enumerate}

    
    \item I have a list of activities that could cause air pollution in the room the hardware is set up in. Can you verify if any of the following occur in the room of the purifier? (\textit{Saya memiliki daftar aktivitas yang dapat menyebabkan polusi udara di ruangan tempat alat dipasang. Bisakah Ibu/Bapak memberitahu apakah terdapat aktivitas berikut yang terjadi di ruangan yang terdapat air purifier?})
    \begin{enumerate}
        \item Gaps in doors and windows (\textit{Celah di pintu dan udara})
        \item Doors/windows are always open (\textit{Pintu/jendela yang selalu terbuka})
        \item Close to kitchen (\textit{Dekat dengan dapur})
        \item Smoking/vaping in or near the room (\textit{Merokok/vaping di dalam/dekat ruangan})
        \item Whether the room has doors and is connected to another room (\textit{Ruangan memiliki pintu yang terhubung dengan ruangan lain})
        \item Garbage burning in the neighborhood (\textit{Pembakaran sampah di lingkungan sekitar})
        \item Candles/Incense (\textit{Difuser/Pelembab udara/humidifier})
        \item None
    \end{enumerate}

\end{enumerate}


\subsection*{Endline Survey Instruments}

Study participants were asked the following survey questions during an endline survey, after the air quality monitoring hardware had been installed for two to three months.\footnote{These questions were part of a longer survey.} Bahasa Indonesian translations are in parentheses.

\begin{enumerate}
    \item In the last 24 hours, did someone smoke inside your home?	(\textit{Dalam 24 jam terakhir, apakah seseorang merokok di dalam rumah Anda?})
    \begin{enumerate}
        \item Yes (\textit{Ya})
        \item No (\textit{Tidak})
         \item Does not know (\textit{Tidak tahu})
        \item Refuse to answer (\textit{Menolak menjawab})
    \end{enumerate}
    
\end{enumerate}

\subsection*{Whatsapp Survey}

Study participants were also asked high frequency questions via Whatsapp during the course of the study. In particular, they were asked about their window/door opening activities:
\begin{enumerate}
    \item In the last hour, was the door to the room where the air quality equipment is stored open to the outside? (\textit{Dalam 1 jam terakhir, apakah pintu di ruangan tempat menyimpan alat kualitas udara terbuka ke luar? (jika ada pintu yang terbuka ke luar ruangan)})
    \begin{enumerate}
        \item Yes (\textit{Ya})
        \item No (\textit{Tidak})
        \item There are no outside doors (\textit{Tiada pintu ke luar})
    \end{enumerate}
        \item In the last hour, was the window in the room where the air quality equipment is stored open to the outside? (\textit{Dalam 1 jam terakhir, apakah jendela di ruangan tempat menyimpan alat kualitas udara terbuka ke luar? (jika ada jendela yang terbuka ke luar ruangan)})
    \begin{enumerate}
        \item Yes (\textit{Ya})
        \item No (\textit{Tidak})
        \item There are no outside windows (\textit{Tak ada jendela luar})
    \end{enumerate}
\end{enumerate}

\clearpage


\newpage
\pagebreak 
\subsubsection*{Envelope Tightening: Back-of-the-envelope.}

Following \citep{chen_review_2011}, in steady-state, indoor PM$_{2.5}$ concentrations $C_{in}$ can be expressed as a function of outdoor PM$_{2.5}$ concentrations $C_{out}$, the air exchange rate due to infiltration \textit{a}, the penetration factor \textit{P}, the particle deposition rate \textit{k} and the indoor-source strength by volume \textit{S}:
\begin{align*}
    C_{in} &= \frac{a P C_{out} +S}{a+k}.
\end{align*}
The coefficient $\beta$ estimated in Eq. (1)  can be shown to correspond to the infiltration factor $F_{inf} = \frac{a P}{a + k}$. For the central back-of-the-envelope calculation, we set $P = 1$ and $k = 0.5 h^{-1}$ and solve for the implied air exchange rate
\begin{align}
    a &= \frac{\hat{F}_{inf} k}{P - \hat{F}_{inf}}
\end{align}
Given $(a, P, k)$, we recover indoor-source strength by volume as the residual
\begin{align*}
    S &= (a + k) C_{in} - a P C_{out}.
\end{align*}
Robust range: we vary the deposition loss rate $k$ over $[0.2, 1.5]~(h^{-1})$ and the penetration factor $P$ over $[0.8, 1.0]$. We evaluate the four corner combinations $(k_{min}, P_{min}), (k_{min}, P_{max}), (k_{max}, P_{min}), (k_{max}, P_{max})$; for each we solve for the implied $a=\frac{\hat F_{inf}k}{P-\hat F_{inf}}$ and recompute $S=(C_{in}-\hat F_{inf}C_{out})(a+k)$.
\\

This model can be leveraged in a simple back-of-the-envelope calculation to assess the impact of interventions that reduce air exchange (targeting \textit{a}) or reduce particle penetration of incoming air (targeting \textit{P}) on indoor PM$_{2.5}$ concentrations. Table \ref{SI_BoE_table} lists the data inputs and external parameters used as well as the main results, split by income and smoking status.
\\

The back-of-the-envelope exercise suggests that reducing particle penetration (lowering $P$) mechanically decreases indoor PM$_{2.5}$ in all households—under the convention that $P$ applies only to outdoor-to-indoor particle entry—whereas weatherization (reducing $a$) can have ambiguous effects because it also limits dilution of indoor emissions. In our central calibration, weatherization increases indoor PM$_{2.5}$ in low-income households, and this conclusion is robust over the parameter range considered. In contrast, weatherization decreases concentrations in high-income non-smoking households throughout the robust range; for high-income smoking households, the effect remains small and can change sign over the robust range. Precisely quantifying the impacts of interventions that alter $a$ and $P$ would require stronger empirical identification of these parameters, which is beyond the scope of this paper.

\begin{table}[!h]
    \caption{Back-of-the-envelope calculations of the effect of weatherization (reducing $a$) and reduced particle penetration of incoming air (reducing $P$) on indoor PM$_{2.5}$ concentrations.}
    \label{SI_BoE_table}

\begin{adjustbox}{width=0.95\linewidth}   % was: max width=\linewidth
  % -----------------------------------------------------------------------
% Auto-generated by table_boe.R -- do not edit by hand
% -----------------------------------------------------------------------

\begin{tabular}{lcccccc}
\toprule
%---- Panel A -------------------------------------------------------
\multicolumn{7}{l}{\textbf{Panel A. Data inputs and derived parameters}} \\[2pt]
  & $C_{in}$
  & $C_{out}$
  & $\hat{F}_{inf}$
  & $a$
  & $S$
  \\
  & ($\mu$g\,m$^{-3}$)
  & ($\mu$g\,m$^{-3}$)
  &
  & (h$^{-1}$)
  & ($\mu$g\,m$^{-3}$\,h$^{-1}$)
  \\
\cmidrule(lr){2-3}\cmidrule(lr){4-4}\cmidrule(lr){5-6}
Income Bin 1, Non-smoking & 48.9 & 39.6 & 0.777 & 1.7 & 40.7 \\
Income Bin 1, Smoking & 57.2 & 40.2 & 0.777 & 1.7 & 58.4 \\
Income Bin 4, Non-smoking & 28.1 & 39.6 & 0.514 & 0.5 & 8.0 \\
Income Bin 4, Smoking & 36.8 & 38.3 & 0.514 & 0.5 & 17.6 \\
\midrule
%---- Panel B -------------------------------------------------------
\multicolumn{7}{l}{\textbf{Panel B. Counterfactual $\Delta C_{in}$ ($\mu$g\,m$^{-3}$)}} \\[2pt]
  &
  \multicolumn{2}{c}{Weatherization ($a \times 0.8$)}
  &
  \multicolumn{2}{c}{Ventilation ($P \times 0.5$)}
  & \\
\cmidrule(lr){2-3}\cmidrule(lr){4-5}
  & Central
  & Robust.~range\textsuperscript{(a)}
  & Central
  & Direction
  & \\
\cmidrule(lr){2-3}\cmidrule(lr){4-5}
Income Bin 1, Non-smoking & $+1.7$ & [$-17.6$,\;$+16.0$]\textsuperscript{\dag} & $-15.4$ & Always $\downarrow$ \\
Income Bin 1, Smoking & $+3.1$ & [$-24.5$,\;$+23.6$]\textsuperscript{\dag} & $-15.6$ & Always $\downarrow$ \\
Income Bin 4, Non-smoking & $-1.3$ & [$-9.1$,\;$+4.4$] & $-10.2$ & Always $\downarrow$ \\
Income Bin 4, Smoking & $-0.2$ & [$-17.3$,\;$+12.5$] & $-9.8$ & Always $\downarrow$ \\
\bottomrule
\end{tabular}

%---- Notes (minipage so footnotes stay attached to table) ----------
\begin{minipage}{\linewidth}
\footnotesize
\textit{Notes.}
$C_{in}$ and $C_{out}$ are mean observed indoor and outdoor PM$_{2.5}$ concentrations
($\mu$g\,m$^{-3}$).
$\hat{F}_{inf}$ is the estimated infiltration factor from regression (see main text).
The air exchange rate $a$ and indoor source strength $S$ are derived from $\hat{F}_{inf}$,
$C_{in}$, $C_{out}$ using Equations~(2)--(4), at the central parameter values
$k = 0.5$ h$^{-1}$ and $P = 1.0$.
For the weatherization scenario, $a$ is reduced by 20\% and
$C_{in}$ is recalculated holding $S$, $k$, and $P$ fixed.
For the ventilation scenario, $P$ is halved and $C_{in}$ is recalculated
holding $S$, $k$, and $a$ fixed; this intervention mechanically reduces
$C_{in}$ regardless of the level of indoor sources.

\smallskip
\textsuperscript{(a)}~Robustness range for the weatherization scenario spans $k \in \{0.3,\,5.0\}$\,h$^{-1}$ and $P \in \{0.7,\,1.0\}$ where feasible.
\smallskip
\textsuperscript{\dag}~For Income Bin 1 ($\hat{F}_{inf} > 0.7$), $P_{\min} = 0.7 < \hat{F}_{inf}$, which would imply a non-positive air exchange rate and is therefore excluded; the robustness range for these groups varies only $k$, holding $P = 1.0$.
\end{minipage}

\end{adjustbox}
\\
  {\footnotesize \textit{Notes:} Panel A reports group means of $C_{in}$ and $C_{out}$ and the estimated infiltration factors $\hat F_{inf}$ by income group. Panel B reports counterfactual changes $\Delta C_{in}$ from (i) weatherization modeled as $a \rightarrow 0.8a$ and (ii) reduced particle penetration of incoming air modeled as $P \rightarrow 0.5P$. \textsuperscript{(a)} Robust range varies $k \in [0.2,1.5]~h^{-1}$ and $P \in [0.8,1.0]$, recomputing $a=\frac{\hat F_{inf}k}{P-\hat F_{inf}}$ and $S=(C_{in}-\hat F_{inf}C_{out})(a+k)$ for each parameter setting as described in the text; the reported interval is the minimum and maximum $\Delta C_{in}$ across these settings.}
\end{table}


\newpage
\subsubsection*{LMG $R^2$ Decomposition.} 
To complement the mean contribution decomposition in the main text, we decompose the $R^2$ of Equation 1 across predictors using the Lindeman-Merenda-Gold (LMG) method \citep{lindeman_introduction_1980}. 
In multiple regressions, the contribution of a predictor to explained variance depends on which other predictors are already in the model, rendering simple sequential decompositions sensitive to arbitrary ordering choices. LMG resolves the attribution problem that arises under predictor correlation by averaging each predictor's marginal $R^2$ contribution over all possible orderings in which predictors could enter the model. This method thus provides a consistent decomposition of explained variance. We implement this using the \texttt{relaimpo} package in \texttt{R} \citep{gromping_relative_2006}. 
The ranking of sources by $R^2$ contribution is consistent with the mean contribution decomposition in the main text: outdoor ambient pollution accounts for the largest share, followed by household smoking and waste burning, with kitchen source, cooking hour, and distance to main road contributing minimally. For pollution spikes, smoking is by far the dominant contributor, while waste burning also contributes substantially.
\begin{table}[!h]
    \centering

    \resizebox{\textwidth}{!}{
    \begin{threeparttable}
    \caption{Source Contribution to Mean Indoor PM$_{2.5}$}     
    \label{tab:source}
    
\begin{tabular}{cccccc}
\toprule
Source & Reg. Estimate & p.value & Mean Value of Source & Mean Contribution & $R^2$ Contribution\\
\midrule
Outdoor Ambient & 0.59 & 0.000 & 39.41 & 56\% & 33\%\\
Smoked (24 Hours) & 21.33 & 0.001 & 0.38 & 19\% & 23\%\\
Kitchen source & -0.28 & 0.968 & 0.31 & 0\% & 0\%\\
Cooking & -2.49 & 0.469 & 0.08 & 0\% & 0\%\\
Distance to Main Road (km) & 3.52 & 0.707 & 0.57 & 5\% & 0\%\\
\addlinespace
Waste Burning (1-2/week) & -1.76 & 0.743 & 0.21 & -1\% & \\
Waste Burning (3+/week) & 46.85 & 0.168 & 0.10 & 12\% & \\
 &  &  &  &  & 22\%\\
\bottomrule
\end{tabular}

    \begin{tablenotes}
      \scriptsize
      \item \textit{Notes:} Regression estimates are from Figure 3a, a fixed-effects model of hourly indoor PM$_{2.5}$ on source variables with hour-of-day and week fixed effects; standard errors are clustered at the household and date-hour level. ``Reg. Estimate" reports the coefficient on each source. "Mean Value of Source" is the sample mean of the corresponding variable (Figure 3b). "Mean Contribution" is the product of the coefficient and the sample mean, expressed as a share of mean indoor PM$_{2.5}$. ``$R^2$ Contribution" reports the proportional contribution of each source to explained variance using the Lindeman-Merenda-Gold (LMG) decomposition from a linear model with hour and week indicators. The LMG decomposition aggregates across levels of categorical variables (e.g., waste burning frequencies), so a single $R^2$ contribution is reported for waste burning overall. 
    \end{tablenotes}
  \end{threeparttable}
    }
\end{table}


\begin{table}[!h]
    \centering

    \resizebox{\textwidth}{!}{
    \begin{threeparttable}
    \caption{Source Contribution to Spikes in Indoor PM$_{2.5}$}   \label{tab:spike_source}
    
\begin{tabular}{cccccc}
\toprule
Source & Reg. Estimate & p.value & Mean Value of Source & Mean Contribution & $R^2$ Contribution\\
\midrule
Outdoor Ambient & 0.000 & 0.008 & 39.41 & 28\% & 2\%\\
Smoking Household & 0.032 & 0.000 & 0.38 & 57\% & 54\%\\
Kitchen source & -0.006 & 0.180 & 0.31 & -9\% & 2\%\\
Cooking & 0.008 & 0.191 & 0.08 & 3\% & 1\%\\
Distance to Main Road (km) & -0.006 & 0.109 & 0.57 & -16\% & 5\%\\
\addlinespace
Waste Burning (1-2/week) & -0.002 & 0.663 & 0.21 & -2\% & \\
Waste Burning (3+/week) & 0.047 & 0.164 & 0.10 & 23\% & \\
 &  &  &  &  & 19\%\\
\bottomrule
\end{tabular}

    \begin{tablenotes}
      \scriptsize
      \item \textit{Notes:} Estimates  from a fixed-effects regression of a binary spike indicator on source variables, with hour-of-day and week fixed effects; standard errors are clustered at the household and date-hour levels. A spike is defined as an hourly PM$_{2.5}$ reading that exceeds the preceding hour by at least one standard deviation.  ``Reg. Estimate" reports the coefficient on each source (from Figure S8). "Mean Value of Source" is the sample mean of the corresponding variable. ``Mean Contribution" is the product of the coefficient and the sample mean, expressed as a share of mean indoor PM$_{2.5}$. ``$R^2$ Contribution" reports the proportional contribution of each source to explained variance using the Lindeman-Merenda-Gold (LMG) decomposition from a linear model with hour and week indicators. The LMG decomposition aggregates across levels of categorical variables (e.g., waste burning frequencies), so a single $R^2$ contribution is reported for waste burning overall. 
    \end{tablenotes}
  \end{threeparttable}
    }
\end{table}


\newpage
\pagebreak 
\section*{Tables}
% \setcounter{table}{0}
% \renewcommand{\thetable}{A\arabic{table}}

\begin{table}[!h]
    \centering

    \resizebox{.7\textwidth}{!}{
    \begin{threeparttable}
    \caption{Summary Statistics of Household Survey}      \label{tab:covariates}
    
\begin{tabular}{lccc}
\toprule
Variable & Control Mean & N & Population Mean\\
\midrule
Female & 0.9 & 283 & \\
 &  &  \vphantom{3} & \\
Age & 41.21 & 283 & \\
 & (6.91) &  & \\
HH Size & 5.13 & 283 & 3.79\\
 & (1.94) &  & (1.31)\\
Number of Children & 2.05 & 283 & \\
 & (1.03) &  & \\
HH Monthly Exp. (USD) & 248.74 & 275 & 655.94\\
 & (117.16) &  & (707.71)\\
HH Monthly Income (USD) & 312.48 & 281 & \\
 & (194.17) &  & \\
HH Head Attended Secondary & 0.88 & 283 & 0.88\\
 &  &  \vphantom{2} & \\
HH Head Attended Tertiary & 0.14 & 283 & 0.3\\
 &  &  \vphantom{1} & \\
Owns AC & 0.55 & 282 & 0.39\\
 & (0.5) &  & \\
LPG Stove & 0.98 & 283 & 0.96\\
 &  &  & \\
Number of Smokers & 0.93 & 283 & 0.73\\
 & (0.85) &  & (0.71)\\
\midrule
\bottomrule
\end{tabular}

    \begin{tablenotes}
      \scriptsize
      \item \textit{Notes:}  Household characteristics and comparisons to the Jakarta population mean. Population means were calculated using data from the National Socioeconomic Survey (SUSENAS) 2023.
    \end{tablenotes}
  \end{threeparttable}
    }
\end{table}


\begin{table}[!h]

    \centering

    \resizebox{.7\textwidth}{!}{
    \begin{threeparttable}
    \caption{Distance to Outdoor Sensor Not Correlated with Household Characteristics}
    \label{tab:sensor_selection}

\begin{tabular}{lc}
   \toprule
   Dependent Variable:           & Distance to Outdoor Sensor (meters)\\  
   Model:                        & (1)\\  
   \midrule
   \emph{Variables}\\
   Indoor PM2.5                  & -0.99\\   
                                 & (2.44)\\   
   Outdoor PM2.5                 & 1.19\\   
                                 & (9.21)\\   
   Income (USD)                  & 0.13\\   
                                 & (0.58)\\   
   Housing Size                  & 42.52\\   
                                 & (53.36)\\   
   AC                            & -198.33\\   
                                 & (197.47)\\   
   HOH Employed                  & 314.69\\   
                                 & (374.62)\\   
   HOH Attended Secondary School & -285.71\\   
                                 & (305.55)\\   
   HOH Attended Tertiary School  & -173.15\\   
                                 & (317.86)\\   
   \midrule
   \emph{Fit statistics}\\
   Dependent variable mean       & 2,337.6\\  
   Observations                  & 283\\  
   \bottomrule
\end{tabular}



    \begin{tablenotes}
      \scriptsize
      \item \textit{Notes:} Regression of households' distance to nearest outdoor sensor on household covariates. IID standard-errors in parentheses. Signif. Codes: $^{***}$: 0.01, $^{**}$: 0.05, $^{*}$: 0.1.
    \end{tablenotes}
  \end{threeparttable}
  }
\end{table}

\begin{table}[!h]
    \centering

    \resizebox{.7\textwidth}{!}{
    \begin{threeparttable}
    \caption{Association Between Reported Waste Burning and Outdoor PM$_{2.5}$}      \label{tab:waste}
    
\begingroup
\centering
\begin{tabular}{lcc}
   \tabularnewline \midrule \midrule
   Dependent Variable: & \multicolumn{2}{c}{Outdoor PM2.5}\\
                            & Closest 3 Sensors & HH Within 2km of Outdoor Sensor \\   
   Model:                   & (1)               & (2)\\  
   \midrule
   \emph{Variables}\\
   Waste Burning (1-2/week) & 0.900             & 0.410\\   
                            & (0.770)           & (0.734)\\   
   Waste Burning (3+/week)  & 3.65$^{***}$      & 5.25$^{***}$\\   
                            & (0.991)           & (1.79)\\   
   Outdoor Temperature      & -0.380            & 0.145\\   
                            & (0.243)           & (0.255)\\   
   Outdoor Humidity         & -0.194$^{***}$    & -0.180$^{***}$\\   
                            & (0.042)           & (0.048)\\   
   \midrule
   \emph{Fixed-effects}\\
   Week FE                  & Yes               & Yes\\  
   Hour FE                  & Yes               & Yes\\  
   \midrule
   \emph{Fit statistics}\\
   Observations             & 823,850           & 423,609\\  
   R$^2$                    & 0.38558           & 0.37695\\  
   Within R$^2$             & 0.00940           & 0.01310\\  
   \midrule \midrule
   \multicolumn{3}{l}{\emph{Clustered (respondent\_id \& date\_hour) standard-errors in parentheses}}\\
   \multicolumn{3}{l}{\emph{Signif. Codes: ***: 0.01, **: 0.05, *: 0.1}}\\
\end{tabular}
\par\endgroup



    \begin{tablenotes}
      \scriptsize
      \item \textit{Notes:}  Each column regresses outdoor PM$_{2.5}$ on household-reported waste burning frequency in the neighborhood over the preceding 7 days, controlling for outdoor temperature and humidity, with hour-of-day and week fixed effects. Column 1 assigns each household the mean PM$_{2.5}$ reading from the three nearest government sensors. Column 2 restricts the sample to households located within 2 km of an outdoor sensor. Standard errors, clustered at the household and date-hour level, are in parentheses. Clustered (respondent and date-hour FE) standard errors are in parentheses. Signif. Codes: $^{***}$: 0.01, $^{**}$: 0.05, $^{*}$: 0.1.
    \end{tablenotes}
  \end{threeparttable}
    }
\end{table}


\clearpage
\section*{Figures}

\begin{figure}[!htb]
    \centering
    \caption{Missing Indoor PM$_{2.5}$  Data} \includegraphics[width=.6\linewidth]{output_clean/figures/fig_appendix_missing_data.png}
    \caption*{\scriptsize Notes: Empirical cumulative distribution function of missing indoor PM$_{2.5}$ data. Each observation represents a household; the x-axis shows the fraction of that household's hourly PM$_{2.5}$ readings that are missing, and the y-axis shows the cumulative share of households at or below that fraction.}
    \label{fig:indoor_missing}
\end{figure}


\begin{figure}[!h]
    \centering
    \caption{Infiltration Rate at Varying Distance to Outdoor Sensor} \includegraphics[width=.6\linewidth]{output_clean/figures/fig_appendix_inf_by_distance.png}
    \caption*{\footnotesize Notes: Each point reports the coefficient on outdoor PM$_{2.5}$ (the infiltration rate) from estimating Equation 1 with hour-of-day and week fixed effects, controlling for hyperlocal source variables, outdoor temperature, and humidity. The sample is progressively restricted to households located within 0.5km, 1km, and 2~km of the nearest outdoor air quality sensor; the rightmost point uses the full sample with no distance restriction. Darker error bars show 90\% confidence intervals, and lighter error bars show 95\% confidence intervals. Standard errors are clustered at the household and date-hour level.}
    \label{fig:dist_sensor}
\end{figure} 


\begin{figure}[!h]
    \centering
    \caption{Infiltration Rate at With Varying Weights} \includegraphics[width=.6\linewidth]{output_clean/figures/fig_appendix_inf_by_weight.png}
    \caption*{\footnotesize Notes: Each column reports the coefficient on outdoor PM$_{2.5}$ (the infiltration rate) from estimating the main infiltration specification. Column 1 is our main specification, which does not have any weights. This implicitly upweights households with more available data. Column 2 weights by the inverse of the household-specific observation rate, upweighting households with more missing data. Column 3 weights by the inverse of the household-by-hour observation rate, upweighting household-hour cells with more missing data. Observation rates are defined as the fraction of non-missing indoor PM$_{2.5}$ readings within the relevant grouping. Darker error bars show 90\% confidence intervals, and lighter error bars show 95\% confidence intervals. Standard errors are clustered at the household and date-hour level.}
    \label{fig:dist_sensor}
\end{figure} 



\begin{figure}[!htb]
    \centering
    \caption{Infiltration Rate with Lagged Outdoor $PM_{2.5}$} \includegraphics[width=.8\linewidth]{output_clean/figures/inf_lags.png}
    \caption*{\scriptsize Notes: Infiltration rate, showing the relationship between changes in outdoor concentration and changes in indoor concentrations. This shows the lags in the previous 11 hours. The model controls for hyperlocal source variables (waste burning frequency, household smoking, kitchen source presence, cooking activity, and distance to main road), outdoor temperature, humidity, and hour-of-day and week fixed effects. Points represent mean estimates; darker error bars show 90\% confidence intervals, and lighter error bars show 95\% confidence intervals. Errors are clustered at the household level. }
    \label{fig:inf_lag}
\end{figure}


\begin{figure}[!htb]
    \centering
    \caption{Robustness: Mean $PM_{2.5}$ Under Various Aggregation Methods} \includegraphics[width=\linewidth]{output_clean/figures/pm_income_robustness.png}
    \caption*{\scriptsize Notes: \textbf{(a)} Mean $PM_{2.5}$ concentrations across income quartiles, aggregated using three different methods to show robustness. \textit{HH, $<2$km from Outdoor Sensor}, calculates a single mean \( PM_{2.5} \) value per household, restricting the sample to households located within 2 km of an outdoor sensor to ensure proximity-based comparability. \textit{HH-Hour-of-Day}, aggregates data by computing the mean \( PM_{2.5} \) concentration for each household at each hour of the day. \textit{HH-Hour-Day} uses mean \( PM_{2.5} \) concentrations at the level of each household-hour-date. These mean $PM_{2.5}$ is then regressed onto the income quartile. Points represent mean estimates; darker error bars show 90\% confidence intervals, and lighter error bars show 95\% confidence intervals. Errors are clustered at the household level.
 }
    \label{fig:pm_income_robust}
\end{figure}


\begin{figure}[!htb]
    \centering
    \caption{CDF of household distance to roads} \includegraphics[width=.6\linewidth]{output_clean/figures/distance_to_road_ecdf.png}
    \caption*{\scriptsize Notes: This figure shows the cumulative distribution of each household's distance to the nearest road, by road type. Motorways are limited-access highways or toll roads, primary roads are major arterial roads connecting districts and cities, secondary roads are collector roads connecting neighborhoods, and tertiary roads are local residential streets.}
    \label{fig:inf_lag}
\end{figure}


\begin{figure}[!htb]
    \centering
    \caption{Source Contributions to PM$_{2.5}$ Spikes} \includegraphics[width=.9\linewidth]{output_clean/figures/fig_appendix_spike_sources.png} 
    \caption*{\scriptsize Notes: Analogous to Figure 3 but with PM$_{2.5}$ spikes as the outcome variable. Spikes are defined as hours in which the hour-to-hour change in indoor PM$_{2.5}$ exceeds one standard deviation of the overall indoor PM$_{2.5}$ distribution. (a) Coefficient estimates from estimating Eq. (1) with the spike indicator as the dependent variable. The dashed green line indicates the sample mean spike probability. (b) Estimated mean contribution of each source to PM$_{2.5}$ spikes, calculated by multiplying each regression coefficient by the sample mean of its corresponding variable and expressing the product as a share of the mean spike probability. Error bars reflect 95\% confidence intervals, computed by scaling the upper and lower bounds of each coefficient's confidence interval by the corresponding sample mean. Standard errors are clustered at the household and date-hour level. }
    \label{fig:spikes}
\end{figure}





   % all SI text here
%%% Add this line AFTER all your figures and tables
\FloatBarrier

% \movie{Type legend for the movie here.}

% \movie{Type legend for the other movie here. Adding longer text to show what happens, to decide on alignment and/or indentations.}

% \movie{A third movie, just for kicks.}

% \dataset{dataset_one.txt}{Type or paste legend here.}

% \dataset{dataset_two.txt}{Type or paste legend here. Adding longer text to show what happens, to decide on alignment and/or indentations for multi-line or paragraph captions.}



% \vspace{1cm}
% \bibsplit[8]
% %Use \bibsplit to split the references from the body of the text. Value "[2]" represents the number of reference in the left column (Note: Please avoid single column figures & tables on this page.)

% Bibliography
\bibliography{JakartaPM}
% \bibliography{pnas-sample}


\end{adjustwidth}
\end{document}
